\section{theory setup}\label{sec:theory-setup}

\subsection{spacetime and boundary geometry}\label{subsec:spacetime-boundary-geometry}

we parameterize global AdS$\displaystyle{_{3}}$ by the standard coordinates $\displaystyle{(t,r,\phi)}$, with ranges $\displaystyle{t\in(-\infty,+\infty)}$ and $\displaystyle{r\in[0,\infty)}$, and angular periodicity $\displaystyle{\phi \sim \phi+2\pi}$. the background metric is

\begin{align}
\mathrm{d}s^{2} & =g^{(0)}_{\mu \nu}\mathrm{d}x^{\mu}\mathrm{d}x^{\nu} \\
 & =-(1+r^{2})\mathrm{d}t^{2}+\dfrac{\mathrm{d}r^{2}}{1+r^{2}}+r^{2}\mathrm{d}\phi ^{2}
\end{align}

we set the AdS radius to unity, $\displaystyle{\ell_{\text{AdS}}=1}$. the regulated spacetime $\displaystyle{M}$ is bounded by two Cauchy surfaces $\displaystyle{\Sigma_{i}}$ and $\displaystyle{\Sigma _{f}}$ and by the timelike cylinder $\displaystyle{\Gamma _{r_{\infty}}}$ at $\displaystyle{r=r_{\infty}}$. the future-pointing unit normal to $\displaystyle{\Sigma}$ is $\displaystyle{\tau ^{\mu}}$, and the outward unit normal to $\displaystyle{\Gamma _{r_{\infty}}}$ is $\displaystyle{n^{\mu}}$. we write

\begin{align}
\sigma _{\mu \nu} & =g_{\mu \nu}+\tau _{\mu}\tau _{\nu}, & \gamma _{\mu \nu} & =g_{\mu \nu}-n_{\mu}n_{\nu}, & q_{\mu \nu} & =g_{\mu \nu}+\tau _{\mu}\tau _{\nu}-n_{\mu}n_{\nu}
\end{align}

for the induced metrics on $\displaystyle{\Sigma,\Gamma _{r_{\infty}}}$ and the corner $\displaystyle{\partial \Sigma}$, respectively. the extrinsic curvature of $\displaystyle{\Gamma _{r_{\infty}}}$ is

\begin{align}
K_{\mu \nu} & =\gamma _{\mu}^{~\rho}\gamma _{\nu}^{~\sigma}\nabla _{\rho}n_{\sigma}, & K=\gamma ^{\mu \nu}K_{\mu \nu}
\end{align}

we restrict $\xi ^{\mu}$ to be an exact Killing vector of the background metric $\displaystyle{g^{(0)}}$. every such vector belongs to the allowed Brown-Henneaux asymptotic-symmetry class. in the coordinates above, its components behave as

\begin{align}
\xi ^{t} & =\mathcal{O}(1), & \xi ^{r} & =\mathcal{O}(r), & \xi ^{\phi} & =\mathcal{O}(1)
\end{align}

\subsection{full theory and asymptotic conditions}\label{subsec:full-theory-asymptotic-conditions}

the renormalized action is

\begin{align}
S[g,\Psi] & =\dfrac{1}{2\kappa ^{2}} \int _{M}\mathrm{d}^{3}x\sqrt{ -g }(R+2) +\dfrac{1}{\kappa ^{2}}\int _{\Gamma} \mathrm{d}^{2}x\sqrt{ -\gamma }(K-1) +S_{\mathrm{m}}[g,\Psi]
\end{align}

we use the normalization $\displaystyle{\kappa ^{2}=8\pi G}$. the metric obeys the usual Brown-Henneaux boundary conditions. writing $\displaystyle{\Delta g=g-g^{(0)}}$, these conditions are

\begin{align}
\Delta g_{tt} & =\mathcal{O}(r^{0}), & \Delta g_{t\phi} & =\mathcal{O}(r^{0}), & \Delta g_{\phi \phi} & =\mathcal{O}(r^{0}) \\
\Delta g_{tr} & =\mathcal{O}(r^{-3}), & \Delta g_{r\phi} & =\mathcal{O}(r^{-3}), & \Delta g_{rr} & =\mathcal{O}(r^{-4})
\end{align}

allowed variations obey the same falloffs.

we define the Hilbert stress tensor and the matter pre-symplectic potential current by the off-shell variation

\begin{align}
\delta S_{\mathrm{m}} & =\dfrac{1}{2}\int _{M}\mathrm{d}^{3}x\sqrt{ -g }T^{\mu \nu}\delta g_{\mu \nu}+\int _{M}\mathrm{d}^{3}x\sqrt{ -g }E_{A}\delta \Psi^{A} \\
 & +\int _{\Sigma _{f}-\Sigma _{i}}\mathrm{d}^{2}x\sqrt{ \sigma }\tau _{\mu}\theta _{\mathrm{m}} ^{\mu}[\delta \Psi,\delta g]-\int _{\Gamma}\mathrm{d}^{2}x\sqrt{ -\gamma }n_{\mu}\theta _{\mathrm{m}} ^{\mu}[\delta \Psi,\delta g]
\end{align}

the matter sector is assumed to satisfy the following conditions:

\begin{enumerate}
\def\labelenumi{\arabic{enumi}.}
\tightlist
\item
  it is diffeomorphism covariant and has a well-defined Hilbert stress tensor.
\item
  the matter sector carries no negative powers of $\displaystyle{\kappa}$, so its leading Hilbert stress tensor satisfies\begin{align}T_{\mu \nu} & =\mathcal{O}(\kappa^{0})\end{align}its conservation on the leading matter equation follows from diffeomorphism covariance, as shown in Appendix~\ref{app:matter-diffeomorphism-current}.
\item
  the non-normalizable matter sources are fixed, and every allowed variation satisfies\begin{align}\left.\delta S_{\mathrm{m}}\right|_{\Gamma} &:=-\int _{\Gamma}\mathrm{d}^{2}x\sqrt{-\gamma} n_{\mu}\theta _{\mathrm{m}}^{\mu}[\delta\Psi,\delta g]=0.\end{align}
\item
  the matter fields have compact support or sufficient normalizable decay, and their diffeomorphism current has no independent asymptotic surface charge. in particular, there is no unscreened gauge charge at infinity.
\end{enumerate}

Appendix~\ref{app:matter-diffeomorphism-current} shows explicitly that diffeomorphism covariance decomposes the matter Noether current into the stress current $\displaystyle{\xi _{\nu}T^{\mu \nu}}$, terms proportional to $\displaystyle{E_{A}}$, and the divergence of the superpotential $\displaystyle{\nabla _{\nu}^{(0)}U^{\mu \nu}_{\xi}}$, where $\displaystyle{U_{\xi}^{\mu \nu}}$ is antisymmetric. imposing the leading matter equation and the asymptotic conditions above then removes the last two contributions: the $\displaystyle{E_{A}}$ terms vanish, and the superpotential has zero integral over $\displaystyle{\partial\Sigma}$.

\subsection{perturbative action}\label{subsec:perturbative-action}

we expand the metric and matter fields as

\begin{align}
g_{\mu \nu} & =g^{(0)}_{\mu \nu}+\kappa h_{\mu \nu}+\kappa ^{2}k _{\mu \nu}+\kappa ^{3}p_{\mu \nu}+\mathcal{O}(\kappa ^{4}) \\
\Psi^{A} & =\Phi ^{A}+\kappa \Xi^{A}+\mathcal{O}(\kappa ^{2})
\end{align}

we impose the metric falloffs order by order in $\displaystyle{\kappa}$, so each of $\displaystyle{h_{\mu \nu}}$, $\displaystyle{k_{\mu \nu}}$, and $\displaystyle{p_{\mu \nu}}$, together with their allowed variations, obeys the same componentwise Brown-Henneaux bounds as $\displaystyle{\Delta g_{\mu \nu}}$.

the bulk action therefore has the ordered form

\begin{align}
S_{\text{bulk}} & =\dfrac{1}{\kappa ^{2}}S^{[-2]}+\dfrac{1}{\kappa}S^{[-1]}+S^{[0]}+\kappa S^{[1]}+\mathcal{O}(\kappa ^{2})
\end{align}

where

\begin{align}
S^{[-2]} & =\dfrac{1}{2}\int _{M}\mathrm{d}^{3}x\sqrt{ -g^{(0)} }\mathcal{L}_{g}^{(0)} \\
S^{[-1]} & =\dfrac{1}{2}\int _{M}\mathrm{d}^{3}x\sqrt{ -g^{(0)} }\mathcal{L}_{g}^{(1)}[h] \\
S^{[0]} & =\dfrac{1}{2}\int _{M}\mathrm{d}^{3}x\sqrt{ -g^{(0)} }(\mathcal{L}_{g}^{(1)}[k]+\mathcal{L}_{g}^{(2)}[h,h])+S_{\mathrm{m}}^{(0)}[g^{(0)},\Phi] \\
S^{[1]} & =\dfrac{1}{2}\int _{M}\mathrm{d}^{3}x\sqrt{ -g^{(0)} }(\mathcal{L}_{g}^{(1)}[p]+2\mathcal{L}_{g}^{(2)}[h,k]+\mathcal{L}_{g}^{(3)}[h,h,h]) \\
 & +\dfrac{1}{2}\int _{M}\mathrm{d}^{3}x\sqrt{ -g^{(0)} }h_{\mu \nu}T^{\mu \nu}+S_{\mathrm{m}}^{(1)}[g^{(0)};\Phi,\Xi]
\end{align}

for our purposes, only $\displaystyle{S^{[-2]},S^{[-1]},S^{[0]}}$, and terms $\displaystyle{\mathcal{L}_{g}^{(2)}[h,k]}$, $\displaystyle{\mathcal{L}_{g}^{(3)}[h,h,h]}$, $\displaystyle{h_{\mu \nu}T^{\mu \nu}}$ in the $\displaystyle{S^{[1]}}$ level are needed.

the timelike-boundary action is expanded in parallel,

\begin{align}
S_{\Gamma}[g] & :=\dfrac{1}{\kappa ^{2}}\int _{\Gamma}\mathrm{d}^{2}x\sqrt{ -\gamma }(K-1) \\
 & =\dfrac{1}{\kappa ^{2}}S_{\Gamma}^{[-2]} +\dfrac{1}{\kappa}S_{\Gamma}^{[-1]}[h] +S_{\Gamma}^{[0]}[h,k] +\kappa S_{\Gamma}^{[1]}[h,k,p] +\mathcal{O}(\kappa ^{2})
\end{align}

Appendix~\ref{app:perturbative-quantities} gives the explicit expressions. the perturbative action $\displaystyle{S=S_{\text{bulk}}+S_{\Gamma}}$ is therefore organized order by order in $\displaystyle{\kappa}$. terms beyond $\displaystyle{\mathcal{O}(\kappa ^{1})}$ are not needed to derive the second-order metric equation used later.

\subsection{goal}\label{subsec:goal}

the standard Brown-York construction, supplemented by holographic renormalization, defines the renormalized boundary stress tensor by varying the renormalized action with respect to the induced boundary metric. in our normalization it is

\begin{align}
\mathcal{T}_{ab} & =\dfrac{1}{\kappa ^{2}}\left(K_{ab}-K\gamma _{ab}+\gamma _{ab}\right)
\end{align}

and the corresponding boundary charge is

\begin{align}
H_{\xi,T}[g] & =-\lim _{r_{\infty}\to\infty}\int _{\partial\Sigma}\mathrm{d}\phi\sqrt{q}\,\tau ^{a}\xi ^{b}\mathcal{T}_{ab}.
\end{align}

we use the \href{https://arxiv.org/abs/1906.08616}{Harlow-Wu covariant phase-space formalism} to construct the bulk Noether charge associated with the diffeomorphism generated by the exact background Killing vector $\displaystyle{\xi}$. our goal is to show that this bulk Noether charge, together with the timelike-boundary and corner contributions fixed by the variational principle, agrees with the Brown-York charge order by order in $\displaystyle{\kappa}$.
